% !TeX program = lualatex
% MIT sample by JetPlugins; tikz-flora copyright its authors, MIT.
\documentclass[11pt,a4paper,twoside]{article}
\usepackage[inner=22mm,outer=57mm,top=39mm,bottom=32mm,
  headheight=14pt,headsep=18pt]{geometry}
% LuaLaTeX's bundled Latin Modern fonts include regular, bold, and italic.
\usepackage{microtype}
\usepackage{fancyhdr}
\usepackage{tikzflora}
\setlength{\parindent}{0pt}
\setlength{\parskip}{9pt}
\pagestyle{fancy}
\fancyhf{}
\fancyhead[LE,RO]{\small\textsc{A season in the garden}}
\fancyfoot[LE,RO]{\small\thepage}
\renewcommand{\headrulewidth}{0pt}

% A fixed seed is repeatable; each page grows a different composition.
% split mirrors the flora into the outer margin of two-sided pages.
% Use seed=random for a new garden each build, or change the number.
\florasetup{layout=split,seed=7,colourway=spring,density=0.8,
  visitors=0.35,giltdpi=300}
\title{The Garden Notebook}
\author{Alex Morgan}
\date{A season of small observations}
\begin{document}

% The title page has its own framing layout and a wider text block.
\newgeometry{left=32mm,right=32mm,top=50mm,bottom=40mm}
\florathispage{layout=flank,seed=17,density=0.9,
  ornaments={sun_crest=4}}
\begin{floratitle}
\maketitle
\thispagestyle{empty}
\end{floratitle}
\vspace{45mm}
\begin{center}
\begin{minipage}{0.68\textwidth}
\centering\itshape
A garden is never quite the same place twice.\par
Leave a little room for what arrives next.
\end{minipage}
\end{center}
\clearpage
\restoregeometry

\section*{At the beginning of spring}
The first green appears close to the ground, where yesterday there was only
soil and last year's leaves. A few shoots lean towards the path. Above them,
a branch carries buds so small that they are easiest to see against the sky.

This notebook begins with those small changes. There is no need to name
every plant or explain every pattern. Write down the color of a stem, the
shape of a leaf, or the place where the afternoon light first reaches the bed.

\subsection*{A practice of looking}
Choose one corner of the garden and return to it each week. Keep the same
point of view, but let the details change. One day the important thing may
be a flower; another day it may be the empty space around it.

\begin{quote}
\itshape Notice first. Name things later. Let a question stay open long
enough for the garden to offer a different answer.
\end{quote}

% Flora also reaches into generous gaps left between blocks of text.
\vspace{28mm}
\subsection*{What to record}
\begin{description}
\item[Light] Where it falls, how long it stays, and what grows towards it.
\item[Weather] Rain on the path, wind at the gate, the warmth of a wall.
\item[Visitors] A bee at a blossom, a bird in the hedge, a trail across a leaf.
\end{description}
\clearpage

% This page alone changes season, plant mix, and ink palette.
% The even and odd body pages demonstrate both outer margins.
\florathispage{colourway=autumn,density=0.7,
  mix={lily=1.4,umbel=1.5,bird=0.5,butterfly=0.4},
  palette={leaves=earth}}
\section*{When the season turns}
By autumn the garden has changed its handwriting. Seed heads replace soft
petals, shadows reach farther across the path, and the leaves gather warm
colors along their edges. The same corner now tells a different story.

Read an earlier entry before adding a new one. Something that seemed
unimportant in spring may explain the shape of the garden now. A stem has
bent, a branch has thickened, and a bare patch has become a place to shelter.

\subsection*{Carry something forward}
Keep one detail for the months ahead: the last flower beside the gate, the
sound of rain in an empty pot, or the brief weight of a bird before it flies.
There is room here for an observation that does not yet belong anywhere else.

\vspace{18mm}
{\small\itshape The next page is an invitation to begin again.}
\end{document}
